\documentclass[11pt,a4paper]{article}
\usepackage{brainteaser}

\begin{document}
\brainteasertitle{2}{Suites de pile ou face: jeu de Penney}

\begin{statementbox}
J1 choisit une suite de trois résultats parmi les huit suites possibles de
pile ou face. J2 voit ce choix, puis choisit une autre suite de longueur trois.

On lance ensuite une pièce équilibrée jusqu'à ce que l'une des deux suites
apparaisse dans trois lancers consécutifs. Les fenêtres peuvent se chevaucher.
J1 choisit $PPP$.

Quelle suite J2 doit-il choisir pour maximiser sa probabilité de gagner, et
quelle est cette probabilité?
\end{statementbox}

\begin{hintsbox}
\begin{enumerate}
  \item Chaque motif a probabilité $1/8$ dans un bloc fixé. Cela rend-il les
  temps de première apparition symétriques?
  \item Tester le choix $FPP$.
  \item Considérer la première apparition de $PPP$. Si elle n'a pas lieu aux
  trois premiers lancers, quel symbole peut précéder ses deux derniers $P$?
  \item Trouver un événement de probabilité $1/8$ sur lequel J2 est obligé de
  perdre, quel que soit son choix.
\end{enumerate}
\end{hintsbox}

\correction

J2 choisit
\[
\boxed{FPP}.
\]

Supposons que $PPP$ apparaisse pour la première fois après le troisième
lancer. Le symbole précédant ses deux derniers $P$ ne peut pas être $P$, sinon
$PPP$ serait apparue un lancer plus tôt. Il est donc nécessairement $F$, ce
qui signifie que $FPP$ vient déjà d'apparaître.

J1 ne peut donc gagner que si les trois premiers lancers sont directement
$PPP$. Cet événement a pour probabilité $1/8$. Par conséquent,
\[
\Pp(\text{victoire de J2})=1-\frac18=\boxed{\frac78}.
\]

Cette valeur est optimale: lorsque les trois premiers lancers sont $PPP$, J1
gagne avant tout motif différent. Tout choix de J2 perd donc avec probabilité
au moins $1/8$, et $FPP$ atteint cette borne.

\begin{keyideabox}
Si J1 choisit le motif $abc$, la meilleure réponse de J2 est
\[
\overline b\,a\,b,
\]
où $\overline b$ désigne le symbole opposé à $b$. L'automatisme est d'étudier
les chevauchements entre suffixes et préfixes, pas seulement la probabilité
d'un bloc isolé.
\end{keyideabox}

\begin{resultbox}
Face à $PPP$, le choix optimal est $FPP$ et la probabilité maximale de victoire
de J2 vaut $7/8$.
\end{resultbox}

\end{document}
